\documentclass[10pt]{article}

\usepackage[margin=0.78in]{geometry}
\usepackage{amsmath,amssymb,booktabs,tabularx,microtype}
\usepackage[hidelinks]{hyperref}
\usepackage{xcolor}
\setlength{\parindent}{0pt}
\setlength{\parskip}{5pt}
\renewcommand{\arraystretch}{1.12}

\newcounter{proposition}
\newenvironment{proposition}[1][]{%
  \refstepcounter{proposition}\par\medskip\noindent
  \textbf{Proposition \theproposition\ifx\relax#1\relax\else\ (#1)\fi.}\ \itshape
}{\par\medskip}
\newcounter{assumption}
\newenvironment{assumption}[1][]{%
  \refstepcounter{assumption}\par\medskip\noindent
  \textbf{Assumption \theassumption\ifx\relax#1\relax\else\ (#1)\fi.}\ \itshape
}{\par\medskip}
\newenvironment{proof}{\par\noindent\textit{Proof.}\ }{\hfill$\square$\par}
\newcommand{\E}{\mathbb{E}}
\newcommand{\Var}{\operatorname{Var}}
\newcommand{\Ind}{\mathbb{1}}
\newcommand{\R}{\mathcal{R}}

\title{Definitive DAPRO with Finite-Sample Expected-Budget Control\\
\large Reconsideration of projection control, constant allocation, and dynamic schedules}
\author{Technical report for the DAPRO manuscript}
\date{August 1, 2026}

\begin{document}
\maketitle

\begin{abstract}
We reconsider the final DAPRO design against two alternatives: the previously
selected projection-reserve method and a constant-continuation policy.  The
new definitive method retains the variance-aligned two-bin allocation model,
but separates 200 fully observed rows into 100 policy-fit and 100 independent
budget-control rows.  A fit-fold-only shared cumulative envelope of twice the
average budget tightens the loss envelope, after which conformal risk control (CRC)
selects a nested cumulative-reach contraction.  The resulting method has a
finite-sample marginal expected-total-budget guarantee under exchangeability,
without a projection-accuracy assumption.  A historical matched 20-seed
confirmation on five real-data settings lowered geometric-mean coverage
variance by 3.9\% relative to CRC constant allocation and won on four settings.
Those runs used the superseded future-path row contraction and must be rerun
before they can validate the corrected causal method.  The
projection-reserve version has nearly identical aggregate performance (4.2\%
lower than CRC constant) but retains an empirically imperfect projection
assumption.  We therefore select the CRC architecture provisionally: its
stronger guarantee remains valid under the shared causal envelope, while the
full LPB empirical matrix is being rerun.  We also prove that the proposed family
$1-\lambda^{\alpha t}$ is independent of $\alpha$ after reparameterization,
and give a Jensen argument explaining why constant allocation is difficult to
beat when policy scores add little target-event information beyond time.
\end{abstract}

\paragraph{Causal-cap provenance.}
The method and guarantee in this document use the shared PAV envelope defined
below.  Result files whose allocator name lacks
\texttt{causal\_shared\_pav\_v1} used a row-specific scale computed from the
complete future score path.  They are retained only as historical comparisons
and are not evidence for online deployability.  The production LPB launcher
uses the suffix \texttt{lpb\_v4\_causal\_pav} so corrected and historical
rows cannot be merged silently.

\section{Final decision}

The public name \textbf{DAPRO} now denotes the following single method.

\begin{center}
\fcolorbox{black}{gray!7}{\begin{minipage}{0.93\linewidth}
\textbf{Definitive CRC-DAPRO.}
Use $N_1=200$ fully observed calibration rows: 100 policy-fit rows and 100
independent budget-control rows.  On the fit fold, define the raw
$\alpha_0=0.10$ target event $A_i$ and weights
\[
W_i=\frac{A_i+0.001}{1.001}.
\]
Fit a two-score-bin continuation table by minimizing
$\sum_i W_i/\pi_i$ under the sequential expected-cost constraint.  Enforce a
trajectory-mixture propensity floor $\epsilon=0.005$.  Using only policy-fit
target masses, learn the shared causal envelope below and intersect each
observed cumulative path with it, so every possible row cost is at most
$K=2b$, where $b$ is the target budget per sample.  Freeze
the resulting policy and the 401-member nested family
\[
R_s(x,t)=\epsilon+s\{R_{\rm cap}(x,t)-\epsilon\},
\qquad s\in\{1,0.9975,\ldots,0\}.
\]
Use the independent control fold and CRC to select the largest feasible $s$.
Deploy that policy unchanged in Phase II.
\end{minipage}}
\end{center}

The implementation exposes this method as
\texttt{DAPRO = DefinitiveCRCDAPRO}.  The assumption-based method remains
available as \texttt{DefinitiveDAPRO}; legacy mean-weight DAPRO remains an
explicit ablation.

\section{Why the variance-aligned objective is correct but incomplete}

For a fixed calibration event $A_i\in\{0,1\}$, let $Z_i$ indicate that the
sequential policy reveals the required endpoint, with inclusion probability
$\pi_i$.  The centered inverse-probability term is
\[
V_i=A_i\left(\frac{Z_i}{\pi_i}-1\right).
\]

\begin{proposition}[Exact fixed-target acquisition variance]
Conditional on the complete trajectories and a predictable frozen policy,
\[
\E[V_i\mid\mathrm{data}]=0,
\qquad
\Var(V_i\mid\mathrm{data})
=A_i\left(\frac1{\pi_i}-1\right).
\]
\end{proposition}

\begin{proof}
$Z_i$ is Bernoulli with mean $\pi_i$, so
$\E[Z_i/\pi_i]=1$ and
$\Var(Z_i/\pi_i)=1/\pi_i-1$.  Use $A_i^2=A_i$.
\end{proof}

Thus $\E[A_i(1/\pi_i-1)]$ is the right mechanistic objective for a
\emph{fixed} target.  Overall mean weight $\E[1/\pi_i]$ is not: it spends
equal effort on non-target rows.  This distinction is visible in the dynamic
schedule experiment.  The complement-power schedule has mean weights as low
as 2--8 in several settings, yet substantially larger $A/\pi$ and roughly
twice the geometric coverage variance of CRC constant allocation.

The fixed-target objective is nevertheless not the variance of the complete
LPB pipeline.  The final calibrated candidate is selected from neighboring
quantiles; data splitting, candidate switching, and test-population variation
remain after the acquisition variance of one anchor is nearly zero.  In the
matched experiment, projection DAPRO often makes the raw-anchor proxy
numerically zero while coverage variance remains comparable to constant.  The
correct conclusion is diminishing returns for this surrogate, not a failure
of Proposition~1.

For a fixed candidate, policy-fit randomness changes variance through policy
regret---the learned propensities may be worse than oracle propensities.  It
does not create a second bias term when the Phase-II policy is non-leaky and
the Horvitz--Thompson identity holds.  For the selected LPB, policy changes can
also move the chosen threshold; that instability is outside Proposition~1.

\section{Policy fitting and the causal shared-envelope cap}

For row $i$, let $L_i=\min(T_i,q_i)$ be its unknown active length and let
$p_{i,t}$ be the predictable conditional continuation probability at
one-based step $t$.  Cumulative reach and expected cost are
\[
R_i(t)=\prod_{u=1}^t p_{i,u},\qquad
\pi_i=R_i(L_i),\qquad
c_i=\sum_{t=1}^{L_i}R_i(t).
\]
The temporal implementation observes $x_{i,t}$ before choosing $p_{i,t}$;
$y_{i,t}$ is revealed only afterward.  The two score bins and their cutpoints
are learned only on the 100-row fit fold.  The target weights are
$W_i=(A_i+0.001)/1.001$, so a target row is about one thousand times more
important than a non-target row while flat zero-target directions remain
regularized.

The cumulative-logit budget correction and trajectory floor are also fit on
the policy-fit fold.  The floor is included inside cost fitting, not added
afterward.  Hence $R_i(t)\ge\epsilon=0.005$, every terminal weight is at most
200, and a reported floored mean weight above $10^5$ is impossible.

The raw fitted policy can have a distribution-free cost bound equal to the
full horizon $T_{\max}=200$.  That makes CRC unnecessarily conservative.  Set
$K=2b$, pool the policy-fit target masses as
$\bar a_t=n_f^{-1}\sum_{i\in\mathcal I_f}a_{i,t}$, and solve
\[
\min_e\sum_{t=1}^{T_{\max}}\frac{\bar a_t}{e_t}
\quad\text{subject to}\quad
1\ge e_1\ge\cdots\ge e_{T_{\max}}\ge\epsilon,
\quad \sum_te_t\le K.
\]
Antitonic PAV followed by one scalar dual search gives the exact shared
envelope $e^*$.  At a reached turn $t$, set
\[
R_{i,\rm cap}(t)=\min\{R_{i,\rm fit}(t),e_t^*\}.
\]
This online map uses only a frozen fit-fold envelope and prefixes observed
through turn $t$.  It never uses $T_i$, a Phase-II label, or a future prefix.
Computing a single row-specific contraction from the full future reach path
would not have this predictability property and is therefore not used.

\begin{proposition}[Label-free support bound]
If $\epsilon q(x)\le K$, then for every possible event time $T$,
\[
\sum_{t=1}^{\min(T,q(x))}R_{\rm cap}(x,t)\le K.
\]
Every further affine contraction toward $\epsilon$ has the same bound.
\end{proposition}

\begin{proof}
Pointwise minimum preserves non-increasing cumulative reach and the floor.
Moreover,
$\sum_{t=1}^{q_i}R_{i,\rm cap}(t)\le\sum_{t=1}^{T_{\max}}e_t^*\le K$.
Every shorter event prefix therefore also costs at most $K$, and any further
affine contraction toward the floor cannot increase the sum.
\end{proof}

The cap is not merely a post-hoc correction.  It reduces the valid candidate
loss envelope from 200 to 40 when $b=20$, and to 20 when $b=10$.  The five-seed
cap screen compared $K=b$ and $K=2b$.  The tighter cap prevented useful
concentration; $K=2b$ was selected and then frozen for the independent
20-seed confirmation.

\section{Finite-sample expected-budget guarantee}

Let the total calibration count be $N=n_f+n_c+m$, with policy-fit size
$n_f=100$, independent control size $n_c=100$, and deployment size $m$.  Let
$B_{\rm tot}=Nb$ and let $C_f$ be the realized fully observed fit-fold cost.
For a control row $Z_i$, let $b_i\in[0,T_{\max}]$ be its fully observed pilot
cost and $c_i(s)\in[0,K]$ the conditional expected cost of fixed candidate
$s$.  Define
\[
\rho=\frac{n_c}{m},\qquad
\ell_i(s)=c_i(s)+\rho b_i,qquad
M=K+\rho T_{\max},\qquad
a=\frac{B_{\rm tot}-C_f}{m}.
\]
Candidates are ordered from aggressive to conservative, so $\ell_i(s)$ is
nested.  CRC selects the first candidate satisfying
\[
\frac{\sum_{i=1}^{n_c}\ell_i(s)+M}{n_c+1}\le a.
\tag{1}
\]

\begin{assumption}[CRC sampling conditions]
Conditional on the policy-fit fold, the control and deployment rows are
exchangeable; the candidate family is fixed before control labels are seen;
losses are nested and bounded by $M$; and at least one candidate satisfies
(1).
\end{assumption}

\begin{proposition}[Marginal expected total budget]
Under Assumption~1, definitive CRC-DAPRO satisfies
\[
\E[C_{\rm total}\mid\mathrm{policy\ fit}]
\le B_{\rm tot}.
\]
The result does not assume projection accuracy.
\end{proposition}

\begin{proof}
The standard finite-sample CRC lemma for a fixed nested family, bounded loss,
and selector (1) gives
$\E[\ell_{\rm new}(\widehat s)\mid\mathrm{fit}]\le a$ by exchangeability.
Since a fresh row has the same pilot-cost distribution as a control row,
\[
\begin{aligned}
\E[C_{\rm total}\mid\mathrm{fit}]
&=C_f+n_c\E[b_{\rm new}]
  +m\E[c_{\rm new}(\widehat s)]\\
&=C_f+m\E[c_{\rm new}(\widehat s)+\rho b_{\rm new}]\\
&\le C_f+ma=B_{\rm tot}.
\end{aligned}
\]
\end{proof}

This is a \emph{marginal} expectation statement over the independent
control/deployment sample.  An individual split may have conditional expected
cost above $B_{\rm tot}$ without contradicting the theorem.  Conversely,
checking only that most splits are below budget is not a proof of the theorem.

\section{What the non-CRC version guarantees}

The assumption-based version uses all 200 Phase-I rows for policy fitting and
reserves $r=1$ expected interaction per deployment row.  If $\widehat c_f$ is
corrected fit-fold cost and $c_d$ is population deployment cost, it enforces
\[
\widehat c_f\le b_2-r,
\qquad
b_2=\frac{B_{\rm tot}-C_1}{m}.
\]
Its expected-budget implication is valid under the precise condition
\[
\E[c_d-\widehat c_f]\le r.
\tag{2}
\]
If the stronger splitwise inequality holds, it gives conditional expected
budget control on that split.  With no reserve, (2) requires nonpositive mean
projection error.  The 10-seed matched screen found that reserve-zero DAPRO
exceeded mean budget on four of five settings (by 0.03--0.47 interactions per
row), so it is not a defensible finite-sample budget method.

Reserve one was empirically safe in mean, but the observed splitwise transfer
inequality held on 85--100\% of the 20 confirmation splits, not all of them.
This supports sensitivity, not a distribution-free guarantee.  We retain the
method as a transparent ablation rather than the public default.

\section{Why constant continuation is strong}

A constant conditional probability $p$ gives $R(t)=p^t$.  It has one fitted
degree of freedom, negligible estimation variance, and automatically spends
more on early endpoints.  The target $A$ corresponds to an early 10\%
quantile, while acquisition is allowed through a broader 56\% prior horizon;
exponential front-loading is therefore already well aligned with the task.

The following idealized result formalizes the value of low capacity.  Let
$H$ be an endpoint/time stratum and let $\Pi$ be terminal propensity under a
possibly score-adaptive policy.

\begin{proposition}[No-signal optimality of stratum-constant reach]
Suppose $A\perp\Pi\mid H$ and policies being compared have the same
$\mu_H=\E[\Pi\mid H]$.  Then
\[
\E\!\left[A\left(\frac1\Pi-1\right)\middle|H\right]
\ge
\Pr(A=1\mid H)\left(\frac1{\mu_H}-1\right),
\]
with equality when $\Pi=\mu_H$ within the stratum.
\end{proposition}

\begin{proof}
Conditional independence factors out $\Pr(A=1\mid H)$.  Jensen's inequality
for the convex function $x\mapsto1/x$ gives
$\E[1/\Pi\mid H]\ge1/\E[\Pi\mid H]$.
\end{proof}

Constant continuation is terminal-propensity constant among rows with the
same active length.  Score adaptation can improve it only if the score
predicts the target event within such groups or if it learns a better
reallocation across time groups.  A 100-row fit fold contains only about
10--20 positive raw-target events in these experiments, so that extra signal
is difficult to estimate.  DAPRO's two bins are a deliberately small attempt
to exploit it; high-capacity projected policies pay more estimation regret.

The result is conditional, not a universal optimality theorem.  It explains
why a simple policy can be near oracle when target signal is weak, and exactly
where DAPRO must find gains.

\section{Dynamic time schedules}

The proposed family has conditional continuation
\[
p_t=1-\lambda^{\alpha t},\qquad \lambda\in[0,1],\ \alpha>0.
\]

\begin{proposition}[Alpha is non-identifiable]
For every fixed $\alpha>0$, the set of schedules above is exactly
$\{1-\beta^t:\beta\in[0,1]\}$.
\end{proposition}

\begin{proof}
Set $\beta=\lambda^\alpha$.  The map
$\lambda\mapsto\beta$ is a bijection of $[0,1]$ for every $\alpha>0$.
\end{proof}

Thus tuning $\lambda$ continuously makes $\alpha=0.5,1,2$ identical.  The
implementation confirmed this: across all tested rows/seeds, the maximum
difference in expected budget was $1.1\times10^{-5}$ and the maximum
difference in mean $A/\pi$ was $2.2\times10^{-7}$.

This family also has a heavy-tail property.  Its cumulative reach is
\[
R(t)=\prod_{u=1}^t(1-\beta^u).
\]
For $0<\beta<1$, the infinite product is positive because
$\sum_u\beta^u<\infty$.  It therefore spends asymptotically linear expected
cost on long trajectories rather than exponentially suppressing them.  This
keeps overall weights small but is poorly matched to early target events.

We also tested the genuinely different family
\[
R_\alpha(t)=p^{t^\alpha},\qquad
p_t=p^{t^\alpha-(t-1)^\alpha}.
\]
$\alpha=1$ is exactly constant continuation; $\alpha>1$ is more front-loaded,
and $\alpha<1$ has a heavier tail.  In the 10-seed five-dataset screen, the
geometric coverage-variance ratios relative to CRC constant were 1.136 for
$\alpha=2$ and 1.177 for $\alpha=0.5$; each won only two datasets.  The
complement-power family had ratio 2.028.  No fixed time-only alternative
consistently surpassed constant continuation.

\section{Matched empirical results}

\textbf{Historical-result warning.}
The tables in this section were generated before the causal-cap correction.
Their DAPRO rows use the legacy future-path contraction and therefore measure
an offline/transductive ablation, not the deployable shared-envelope method
defined above.  They remain useful for diagnosing the objective and CRC
tradeoffs, but all headline comparisons below are provisional until the
versioned \texttt{lpb\_v4\_causal\_pav} matrix has been rerun.

All confirmation results use calibration size 3,000, common data splits and
acquisition uniforms, 20 seeds, and five cached real-data configurations:
toxicity/Detoxify, AutoIF, hallucination/Qwen, and red-team data judged by
Llama-Guard or Qwen.  Budgets are 20 for toxicity, AutoIF, and red-team/Qwen,
and 10 for hallucination and red-team/Llama-Guard.

\begin{table}[t]
\centering
\small
\caption{Coverage variance in squared percentage points and mean expected
budget per calibration row.  P is projection-reserve DAPRO, D is definitive
CRC-DAPRO, and C is CRC constant continuation.}
\label{tab:main}
\begin{tabular}{lrrrrrrr}
\toprule
Dataset & Var P & Var D & Var C & Budget P & Budget D & Budget C & Target\\
\midrule
Toxicity & 0.371 & 0.349 & 0.310 & 19.12 & 19.81 & 18.15 & 20\\
AutoIF & 0.728 & 0.731 & 0.915 & 19.18 & 19.77 & 18.34 & 20\\
Hallucination & 0.635 & 0.694 & 0.727 & 9.43 & 9.98 & 8.49 & 10\\
Red-team/Llama-Guard & 0.814 & 0.871 & 0.907 & 9.18 & 9.28 & 8.02 & 10\\
Red-team/Qwen & 0.521 & 0.480 & 0.482 & 19.09 & 19.85 & 18.13 & 20\\
\bottomrule
\end{tabular}
\end{table}

Relative to CRC constant, the geometric variance ratio is 0.958 for the
projection-reserve version and 0.961 for definitive CRC-DAPRO.  The former
wins on three datasets and the latter on four.  Their aggregate difference is
only 0.3 percentage points.  Definitive CRC-DAPRO loses on toxicity but wins
by 0.5--20\% on the other four point comparisons.  These estimates do not
support universal dominance; they support near-parity with a small average
gain and a stronger theorem.

\begin{table}[t]
\centering
\small
\caption{Mean inverse propensity and fixed raw-target $A/\pi$ diagnostic.}
\label{tab:weights}
\begin{tabular}{lrrrrrr}
\toprule
Dataset & Weight P & Weight D & Weight C & $A/\pi$ P & $A/\pi$ D & $A/\pi$ C\\
\midrule
Toxicity & 25.72 & 19.74 & 113.90 & 0.1107 & 0.1375 & 0.1125\\
AutoIF & 20.72 & 12.40 & 110.62 & 0.0988 & 0.1276 & 0.1076\\
Hallucination & 2.05 & 1.92 & 6.43 & 0.0971 & 0.1045 & 0.1042\\
Red-team/Llama-Guard & 3.62 & 3.08 & 26.51 & 0.0999 & 0.1037 & 0.1076\\
Red-team/Qwen & 9.23 & 6.19 & 40.09 & 0.0995 & 0.1138 & 0.1285\\
\bottomrule
\end{tabular}
\end{table}

CRC-DAPRO has substantially smaller mean weights than CRC constant on all
five datasets.  Its $A/\pi$ diagnostic is better on both red-team settings,
approximately tied on hallucination, and worse on toxicity and AutoIF.  Yet
coverage variance improves on AutoIF.  This is direct evidence that the fixed
raw-target variance term is not a complete ordering statistic for the
selected LPB.

The cap materially improves budget use.  The un-capped CRC candidate paid a
correction near 2.05 interactions per row and had geometric variance ratio
1.29--1.41 versus CRC constant, depending on the fit split.  With $K=2b$, the
correction is 0.467 on budget-20 tasks and 0.269 on budget-10 tasks.  The
selected scale averages 0.75--0.99 across datasets.  The final method focuses
more budget on score-favored rows while retaining the support bound.

The projection method's observed transfer inequality held on 85--100\% of
splits.  Definitive CRC-DAPRO's mean expected cost was below target on all
five settings.  Splitwise exceedance rates ranged from 0 to 45\%, which is
compatible with a marginal guarantee and illustrates why the theorem should
not be described as per-split control.

Runtime is also favorable relative to projection DAPRO.  Depending on the
dataset, definitive CRC-DAPRO averages 0.97--5.65 seconds per method/seed,
versus 1.27--10.25 seconds for projection DAPRO.  Constant allocation remains
much faster at roughly 0.05 seconds.

\section{Recommendation and manuscript claims}

Definitive CRC-DAPRO offers the best balance requested:
\begin{itemize}
\item \textbf{Theoretical correctness:} finite-sample marginal expected-total-
budget control without projection accuracy; predictable IPCW acquisition and
strict positivity are preserved.
\item \textbf{Historical empirical signal:} the superseded cap obtained a
3.9\% geometric variance reduction and four of five point wins versus the
guarantee-matched constant policy; the corrected LPB matrix is pending.
\item \textbf{Simplicity:} two score bins, one fixed row cap, and one scalar
CRC contraction; no per-time projection regression.
\item \textbf{Numerical stability:} $\pi\ge0.005$ caps every weight at 200;
the floor and row cap are included inside expected-cost accounting.
\item \textbf{Reproducibility:} all defaults are fixed, all candidate families
are named, common acquisition randomness is stored, and Linux/PowerShell
runners reproduce the screens and confirmation.
\end{itemize}

The historical headline was:
\begin{quote}
Across five 20-seed real-data settings, row-capped CRC-DAPRO had 3.9\% lower
geometric-mean coverage variance than CRC constant allocation and won four
point comparisons, while providing a finite-sample marginal expected-budget
guarantee without a projection-accuracy assumption.
\end{quote}
This sentence must not be used as the corrected method's publication claim
until those five settings have been rerun with an allocator name containing
\texttt{causal\_shared\_pav\_v1}.  At present, the theorem applies to the
corrected policy and the quoted numbers apply only to the historical policy.

Do not claim that DAPRO universally dominates constant allocation.  Do not
use mean weight as the variance objective.  Do not claim that an individual
CRC split must be under budget.  The constant method should remain the primary
baseline because Proposition~5 explains why it can be essentially optimal
when the learned score has little within-time target signal.

\section{Implementation and reproduction}

The core implementation and tests are:
\begin{verbatim}
src/predictive_bounds/budget_allocators/DAPRO.py
src/predictive_bounds/budget_allocators/risk_controlled_budget.py
src/predictive_bounds/budget_allocators/random_adaptive_optimized_allocator.py
tests/test_definitive_dapro.py
tests/test_risk_controlled_budget.py
tests/test_random_allocator_floor_modes.py
\end{verbatim}

The historical screen and confirmation packages are under
\begin{verbatim}
src/predictive_bounds/experiments/crc_dynamic_schedules/
\end{verbatim}
with Linux and PowerShell runners.  Principal machine-readable outputs are
\begin{verbatim}
outputs/crc_dynamic_schedules/method_summary.csv
outputs/crc_row_caps/method_summary.csv
outputs/crc_rowcap2_confirm/method_summary.csv
outputs/crc_rowcap2_confirm/geometric_variance_comparison.csv
\end{verbatim}

The corrected production defaults are $N_1=200$, fit/control sizes 100/100,
$\alpha_0=0.10$, global regularization $0.001$, two score bins,
$\epsilon=0.005$, row cap $K=2b$, and 401 CRC candidates.  Alternative
schedules, un-capped CRC, larger fit folds, and the projection-reserve method
are ablations rather than hidden switches in the paper's algorithm.

\begin{thebibliography}{2}
\bibitem{dapro-manuscript}
S. Feldman and Y. Romano.
\newblock How Many Iterations to Jailbreak? Dynamic Budget Allocation for
Multi-Turn LLM Evaluation.
\newblock arXiv:2605.06605, version 5, 2026.

\bibitem{crc}
A. N. Angelopoulos, S. Bates, A. Fisch, L. Lei, and T. Schuster.
\newblock Conformal Risk Control.
\newblock International Conference on Learning Representations, 2024.
\end{thebibliography}

\end{document}
